\section{Pruebas y calidad}

La estrategia de pruebas del Mini Banco se apoya en un \emph{harness} propio
escrito en Java puro, sin JUnit ni ninguna otra biblioteca externa. Esta fue una
decision deliberada: el proyecto no agrega dependencias de pruebas para no
inflar el \codigo{pom.xml} ni el empaquetado, y porque las suites necesitan
levantar componentes reales (servidor NIO, sockets TCP, repositorio
concurrente) que un framework no aporta. Todas las pruebas viven en el paquete
\codigo{pruebas}, dentro de \codigo{src/test/java/pruebas/}, y comparten dos
piezas de infraestructura: un corredor unico y un verificador minimo de
aserciones.

\subsection{El harness: TestRunner y MiniTest}

El corredor \codigo{pruebas/TestRunner.java} invoca en secuencia el metodo
estatico \codigo{run()} de cada suite, imprime un encabezado por suite y al
final resume cuantas aserciones pasaron o fallaron. Si hubo al menos un fallo,
termina con \codigo{System.exit(1)}, de modo que un script o un pipeline puede
detectar la regresion por el codigo de salida.

\begin{lstlisting}[language=Java,caption={TestRunner orquesta las suites y propaga el fallo via codigo de salida}]
System.out.println("== CuentaRepositoryTest ==");
CuentaRepositoryTest.run();
// ... el resto de las suites ...
System.out.printf("%nResultado: %d pasadas, %d fallidas%n",
        MiniTest.pasadas(), MiniTest.fallidas());
if (MiniTest.fallidas() > 0) System.exit(1);
\end{lstlisting}

El verificador \codigo{pruebas/MiniTest.java} sustituye a las anotaciones y
\emph{assertions} de un framework con apenas tres metodos: \codigo{check()} para
una condicion booleana y dos sobrecargas de \codigo{eq()} (una para
\codigo{long} y otra para \codigo{Object}) que comparan un valor esperado contra
el real. Lleva un par de contadores estaticos, \codigo{pasadas} y
\codigo{fallidas}, e imprime el mensaje de la asercion que no se cumple en lugar
de lanzar una excepcion, lo que permite que una suite reporte todas sus fallas
en una sola corrida.

\begin{lstlisting}[language=Java,caption={MiniTest: aserciones y conteo en Java puro}]
public static void check(boolean cond, String msg) {
    if (cond) {
        pasadas++;
    } else {
        fallidas++;
        System.out.println("  FALLO: " + msg);
    }
}

public static void eq(long esperado, long real, String msg) {
    check(esperado == real, msg + " (esperado " + esperado + ", real " + real + ")");
}
\end{lstlisting}

\subsection{Como se ejecutan}

Toda la bateria se corre con un solo comando, \codigo{./ejecutar\_tests.sh}, que
primero compila con Maven y luego lanza el corredor con el classpath del jar mas
las clases de prueba.

\begin{lstlisting}[language=none,caption={ejecutar\_tests.sh: compila y corre el harness}]
mvn -q package
java -cp "target/banco.jar:target/test-classes" pruebas.TestRunner
\end{lstlisting}

\begin{nota}
En total las ocho suites suman \textbf{68 aserciones} verdes. El conteo final lo
imprime el propio \codigo{TestRunner} y, mientras todas pasen, el script termina
con codigo de salida 0.
\end{nota}

\subsection{Que cubre cada suite}

Las suites van de la logica de dominio aislada hasta los componentes
distribuidos ejercitados con red real. Cada una se cita por su archivo fuente.

\begin{itemize}
\item \codigo{pruebas/CuentaRepositoryTest.java} (13 aserciones): valida la
transferencia basica en centavos, los casos borde (\codigo{MISMA\_CUENTA},
\codigo{MONTO\_INVALIDO}, \codigo{NO\_ENCONTRADA}, \codigo{SALDO\_INSUFICIENTE})
y, sobre todo, la invariante de conservacion del dinero lanzando 8 hilos con
5000 operaciones cada uno (40 mil transferencias concurrentes) y comprobando que
\codigo{saldoTotalCentavos()} no cambia.

\item \codigo{pruebas/AlumnosLoaderTest.java} (7 aserciones): carga un CSV
temporal con \codigo{AlumnosLoader.cargar()} y verifica el numero de filas, la
numeracion de cuentas base 1, el armado del propietario y la conversion exacta a
centavos (por ejemplo \codigo{36180311} para \codigo{361803.11}).

\item \codigo{pruebas/ReplicacionTest.java} (11 aserciones): comprueba que el
log de transacciones registra en orden y que una replica converge al lider
reaplicando \codigo{desde(0)}, incluido el escenario de recuperacion
incremental (catch-up) en que la replica se quedo en la secuencia 150 y pide
solo de la 151 en adelante.

\item \codigo{pruebas/ReplicacionTcpTest.java} (4 aserciones): la misma
convergencia pero con \emph{sockets reales}. Levanta \codigo{ServidorReplicacion}
en un puerto efimero y un \codigo{ClienteReplicacion} que se conecta por TCP a
\codigo{127.0.0.1}; valida el catch-up inicial y que el \emph{streaming} en vivo
propague nuevas transferencias.

\item \codigo{pruebas/GeneradorCargaTest.java} (5 aserciones): corre el
\codigo{GeneradorCarga} unos 2 segundos contra un nodo HTTP real en proceso
(\codigo{BancoHttp.crear(0)}) y verifica que hubo lecturas y transferencias
exitosas sobre la pila HTTP, que descubrio el numero de cuentas via
\codigo{/stats} y que la invariante de suma se conserva.

\item \codigo{pruebas/PanelTest.java} (8 aserciones): comprueba que
\codigo{/panel} agrega el \codigo{/stats} local, el de un peer vivo y marca como
caido un peer inalcanzable, ademas de confirmar que \codigo{dashboard.html} viaja
en el jar y consume \codigo{/panel} de forma reactiva con \codigo{setInterval}.

\item \codigo{pruebas/AlmacenGcsTest.java} (10 aserciones): prueba el log durable
con un \codigo{Subidor} falso (sin tocar GCS ni la red) que puede fallar las
primeras N veces. Verifica persistencia en orden, reintentos tras fallos
transitorios, el sembrado del contador con objetos existentes, el enganche al
commit via observador y la recuperacion reaplicando el log.

\item \codigo{pruebas/ParserHttpTest.java} (10 aserciones): ejercita el parser
HTTP/1.1 del servidor NIO con \emph{sockets crudos}, en los casos borde que un
cliente de alto nivel no toca: dos requests por \codigo{keep-alive}, pipelining
en un mismo segmento TCP, body partido en dos segmentos, respuesta 405 con
\codigo{Content-Length: 0} que no rompe la conexion, query string ignorada,
header \codigo{Authorization} sin distincion de mayusculas y ausencia de
autorizacion (401).
\end{itemize}

\begin{lstlisting}[language=Java,caption={ParserHttpTest: dos requests en una sola conexion keep-alive sobre socket crudo}]
try (Socket s = new Socket("127.0.0.1", puerto)) {
    OutputStream out = s.getOutputStream();
    enviar(out, "GET /stats HTTP/1.1\r\nHost: x\r\n\r\n");
    MiniTest.eq(200, leerResp(s.getInputStream())[0], "keep-alive: 1a request GET /stats -> 200");
    enviar(out, "GET /stats HTTP/1.1\r\nHost: x\r\n\r\n");
    MiniTest.eq(200, leerResp(s.getInputStream())[0], "keep-alive: 2a request en la MISMA conexion -> 200");
}
\end{lstlisting}

\begin{consejo}
La mezcla de pruebas de dominio puro (\codigo{CuentaRepositoryTest}) con pruebas
de integracion sobre red real (\codigo{ReplicacionTcpTest},
\codigo{GeneradorCargaTest}, \codigo{ParserHttpTest}) da confianza tanto en la
logica como en el transporte, sin pagar el costo de un framework externo.
\end{consejo}
